% ============================================================
% f2012_liaoning_science_q09_flowchart.tex
% 2012 辽宁卷（理科）第 9 题：程序框图
% 来源 PDF：普通高考/2012/2012辽宁理.pdf
% 原页：第 1 页（A3），印面 751，第 9 题题干下方
% 初始 S=4、i=1；判断 i<9，是则 S=2/(2-S)、i=i+1 后循环，
% 否则输出 S 并结束。
% ============================================================
\begin{tikzpicture}[>=Stealth, line width=0.7pt, font=\small,
  every node/.style={inner sep=2pt},
  term/.style={draw, rounded corners=3pt, minimum width=2.0cm,
    minimum height=0.58cm, align=center},
  proc/.style={draw, minimum width=2.6cm, minimum height=0.58cm, align=center},
  io/.style={draw, trapezium, trapezium left angle=70, trapezium right angle=110,
    minimum width=2.4cm, minimum height=0.58cm, align=center},
  dec/.style={draw, diamond, aspect=2, minimum width=2.4cm,
    minimum height=0.9cm, align=center}]
  \node[term] (start) at (0,0) {开始};
  \node[proc] (initS) at (0,-0.92) {$S=4$};
  \node[proc] (initI) at (0,-1.84) {$i=1$};
  \node[dec] (test) at (0,-3.02) {$i<9$};
  \node[proc] (bodyS) at (0,-4.46) {$S=\dfrac{2}{2-S}$};
  \node[proc] (bodyI) at (0,-5.52) {$i=i+1$};
  \node[io] (out) at (3.35,-3.02) {输出 $S$};
  \node[term] (end) at (3.35,-4.22) {结束};

  \draw[->] (start) -- (initS);
  \draw[->] (initS) -- (initI);
  \draw[->] (initI) -- (test);
  \draw[->] (test.south) -- node[right] {是} (bodyS.north);
  \draw[->] (bodyS) -- (bodyI);
  % 循环体回到判断框左侧
  \draw[->] (bodyI.west) -- ++(-2.45,0) |- (test.west);
  % “否”分支输出
  \draw[->] (test.east) -- node[above] {否} (out.west);
  \draw[->] (out) -- (end);
\end{tikzpicture}
